\documentclass[]{article}

\usepackage{amsmath}
\usepackage{amsfonts}
\usepackage{amsthm}
\usepackage{amssymb}
\usepackage{mathrsfs}
\usepackage{stmaryrd}

\newcommand{\Q}{\mathbb{Q}}
\newcommand{\N}{\mathbb{N}}
\newcommand{\Z}{\mathbb{Z}}
\newcommand{\R}{\mathbb{R}}
\newcommand{\Primes}{\mathbb{P}}
\newcommand{\st}{\text{ s.t. }}
\newcommand{\txtand}{\text{ and }}
\newcommand{\txtor}{\text{ or }}
\newcommand{\lxor}{\veebar}

%opening
\title{Notes on Differential Equations}
\author{DSBA Mathematics Refresher 2025}
\date{}

\begin{document}
	
	\maketitle
	
	\begin{abstract}
		
	\end{abstract}
	
	\section{Applications of Differential Equations}
	Differential equations play a significant role in various applications within data science and business. These applications demonstrate the versatility of differential equations in modeling complex, dynamic systems across diverse fields, aiding in decision-making and strategic planning.
	\subsection{Predictive Modeling in Finance}
	\paragraph{Black-Scholes Model}\noindent\\
	This is a partial differential equation used to estimate the price of options and other financial derivatives.
	It helps traders and analysts predict the future behavior of asset prices and manage financial risk.
	\subsection{Economics and Market Dynamics}
	\begin{itemize}
		\item \textbf{Macroeconomic Modeling}\noindent\\
		Differential equations are used to model economic growth, inflation, and interest rates.
		For example, the Solow growth model uses differential equations to describe how capital accumulation and technological progress impact economic growth.
		\item \textbf{Epidemic Models}\\
		In business, especially during crises like pandemics, models such as SIR (Susceptible-Infected-Recovered) are used to predict the spread of disease and its impact on markets and workforce productivity.
	\end{itemize}
	\subsection{Customer Behavior and Marketing}
	\begin{itemize}
		\item \textbf{Diffusion Models}\noindent\\
		Differential equations are used to model the adoption of new products and technologies in markets.
		The Bass Diffusion Model, for instance, predicts the rate at which a new product will be adopted over time, helping companies strategize marketing and sales efforts.
		\item \textbf{Churn Prediction}\noindent\\
		Models based on differential equations can predict customer churn by analyzing the rates of customer acquisition and departure over time, allowing businesses to implement targeted retention strategies.
	\end{itemize}
	\subsection{Supply Chain and Inventory Management}
	\paragraph{Dynamic Systems}\noindent\\
	Differential equations are used to model inventory levels over time, helping businesses optimize ordering strategies and reduce holding costs.
	The equations can account for varying demand rates and lead times, improving supply chain efficiency.
	\subsection{Environmental Analysis}
	\paragraph{Time Series Analysis}\noindent\\
	Differential equations are used to model and predict changes in environmental data, such as air quality, temperature, or pollution levels.
	These models help businesses in energy and agriculture make data-driven decisions.

	\subsection{A Concrete Example: Modeling a Customer Base}

	\noindent \textbf{Example:}
	A subscription company acquires $a = 500$ new customers every month, and loses each month a fraction $b = 10\% = 0.1$ of the customers it currently has.
	Writing $N(t)$ for the number of customers at time $t$ (in months), the \textit{rate of change} of the customer base is
	$$
	\underbrace{\frac{dN}{dt}}_{\text{change in customers}} = \underbrace{a}_{\text{acquired}} - \underbrace{b \, N}_{\text{churned}}
	$$
	This is a first-order differential equation; we will see in the next section how to solve it (with the integrating factor $\text{IF} = e^{bt}$).
	Starting from $N(0) = 2000$ customers, the solution turns out to be:
	$$
	N(t) = \frac{a}{b} + \left(N(0) - \frac{a}{b}\right)e^{-bt} = 5000 - 3000 \, e^{-0.1 t}
	$$
	Note what the model tells the business: the customer base does not grow forever, but converges to the equilibrium $\frac{a}{b} = 5000$ customers, at a speed governed by the churn rate $b$.
	Doubling the acquisition rate $a$ doubles the equilibrium; halving the churn rate $b$ also doubles it.
	
	
	
	\section{First-Order Differential Equations}
	
	A first-order differential equation can often be written in the form:
	$$
	\frac{dy}{dx} + P(x) \cdot y = Q(x)
	$$
	To solve this, we use an integrating factor "$\text{IF}$" which is defined as:
	$$
	\text{IF} = e^{\int P(x) \, dx}
	$$
	Multiplying both sides of the differential equation by $\text{IF}$ gives:
	$$
	\text{IF} \cdot \frac{dy}{dx} + \text{IF} \cdot P(x) \cdot y = \text{IF} \cdot Q(x)
	$$
	The left side of the equation is now the derivative of the product $\text{IF} \cdot y$:
	$$
	\frac{d}{dx}\left[\text{IF} \cdot y\right] = \text{IF} \cdot Q(x)
	$$
	Integrating both sides with respect to $x$:
	$$
	\text{IF} \cdot y = \int \text{IF} \cdot Q(x) \, dx + C
	$$
	Finally, solve for $y(x)$:
	$$
	y(x) = \frac{1}{\text{IF}}\left[\int \text{IF} \cdot Q(x) \, dx + C\right]
	$$

	\noindent \textbf{Example:}
	Solve
	$$
	\frac{dy}{dx} + 2x \, y = x
	$$
	Here $P(x) = 2x$ and $Q(x) = x$, so the integrating factor is
	$$
	\text{IF} = e^{\int 2x \, dx} = e^{x^2}
	$$
	Multiplying the equation by $e^{x^2}$:
	$$
	e^{x^2}\frac{dy}{dx} + 2x \, e^{x^2} y = x \, e^{x^2}
	\qquad \Longleftrightarrow \qquad
	\frac{d}{dx}\left[e^{x^2} y\right] = x \, e^{x^2}
	$$
	Integrating both sides (the right-hand side by substitution $u = x^2$):
	$$
	e^{x^2} y = \frac{1}{2} e^{x^2} + C
	\qquad \Longleftrightarrow \qquad
	y(x) = \frac{1}{2} + C e^{-x^2}, \qquad C \in \R
	$$
	\textit{Check:} $y' = -2Cx e^{-x^2}$, so $y' + 2xy = -2Cx e^{-x^2} + x + 2Cx e^{-x^2} = x$. \checkmark

	\noindent \textbf{Example (the customer base of Section 1):}
	The equation $\frac{dN}{dt} = a - bN$ is rewritten in the standard form as $\frac{dN}{dt} + bN = a$, so $\text{IF} = e^{\int b \, dt} = e^{bt}$ and
	$$
	\frac{d}{dt}\left[e^{bt} N\right] = a \, e^{bt}
	\quad \Longrightarrow \quad
	e^{bt} N = \frac{a}{b} e^{bt} + C
	\quad \Longrightarrow \quad
	N(t) = \frac{a}{b} + C e^{-bt}
	$$

	\paragraph{Initial Conditions}
	
	Consider the same first-order equation with an initial condition $y(x_0) = y_0$.
	After finding the general solution as in the previous section:
	$$
	y(x) = \frac{1}{\text{IF}}\left[\int \text{IF} \cdot Q(x) \, dx + C\right]
	$$
	Apply the initial condition to find $C$:
	$$
	y(x_0) = \frac{1}{\text{IF}(x_0)}\left[\int \text{IF}Q(x) \, dx \mid_{x_0} + C\right] = y_0
	$$
	Solve for $C$ and conclude.

	\noindent \textbf{Example:}
	Take the equation of the previous example, $\frac{dy}{dx} + 2xy = x$, with the initial condition $y(0) = 2$.
	The general solution is $y(x) = \frac{1}{2} + C e^{-x^2}$, so
	$$
	y(0) = \frac{1}{2} + C e^{0} = \frac{1}{2} + C = 2
	\qquad \Longrightarrow \qquad
	C = \frac{3}{2}
	$$
	and the solution to the initial value problem is
	$$
	y(x) = \frac{1}{2} + \frac{3}{2} e^{-x^2}
	$$

	\noindent \textbf{Example (the customer base of Section 1):}
	With $N(t) = \frac{a}{b} + Ce^{-bt}$ and $N(0) = N_0$, we get $\frac{a}{b} + C = N_0$, i.e. $C = N_0 - \frac{a}{b}$, hence
	$$
	N(t) = \frac{a}{b} + \left(N_0 - \frac{a}{b}\right) e^{-bt}
	$$
	With $a = 500$, $b = 0.1$ and $N_0 = 2000$: $\frac{a}{b} = 5000$ and $N(t) = 5000 - 3000\, e^{-0.1t}$, as announced.

	\section{Separable Differential Equations}

	The integrating factor of the previous section requires the equation to be \textbf{linear} in $y$.
	When it is not, another method is often available: \textbf{separation of variables}.

	A first-order equation is said to be \textbf{separable} if it can be written in the form
	$$
	\frac{dy}{dx} = f(x) \cdot g(y)
	$$
	that is, if the right-hand side splits into a part depending only on $x$ and a part depending only on $y$.
	The method is then to gather each variable on its own side of the equation, and to integrate:
	$$
	\frac{dy}{g(y)} = f(x) \, dx
	\qquad \Longrightarrow \qquad
	\int \frac{dy}{g(y)} = \int f(x) \, dx
	$$
	Both integrals are ordinary integrals in a single variable.
	One then solves the resulting relation for $y$, if possible, and a single constant of integration $C$ is enough (the two constants can be merged into one).

	\noindent \textbf{Example:}
	Solve
	$$
	\frac{dy}{dx} = x \, y^2
	$$
	This equation is \textit{not} linear in $y$ (because of the $y^2$), so the integrating factor does not apply.
	It is however separable, with $f(x) = x$ and $g(y) = y^2$.
	Dividing by $y^2$ and integrating:
	$$
	\frac{dy}{y^2} = x \, dx
	\qquad \Longrightarrow \qquad
	\int \frac{dy}{y^2} = \int x \, dx
	\qquad \Longrightarrow \qquad
	-\frac{1}{y} = \frac{x^2}{2} + C, \qquad C \in \R
	$$
	Solving for $y$ (and writing $K = 2C$ to simplify):
	$$
	y(x) = \frac{-2}{x^2 + K}
	$$
	\textit{Check:} $y' = \frac{4x}{(x^2+K)^2}$, and $x y^2 = x \cdot \frac{4}{(x^2+K)^2} = \frac{4x}{(x^2+K)^2}$. \checkmark

	\noindent \textbf{Warning:}
	dividing by $g(y)$ is only legitimate when $g(y) \neq 0$.
	The values of $y$ for which $g(y) = 0$ give \textit{constant} solutions that the method silently drops:
	here $y(x) = 0$ is also a solution of $y' = xy^2$, and it is not of the form found above.
	Always check whether such constant solutions exist.

	\noindent \textbf{Example (the two methods agree):}
	The equation $\frac{dy}{dx} = 5y$ is both linear and separable.
	Separating:
	$$
	\frac{dy}{y} = 5 \, dx
	\qquad \Longrightarrow \qquad
	\ln|y| = 5x + C
	\qquad \Longrightarrow \qquad
	|y| = e^{5x + C} = e^{C} e^{5x}
	$$
	so $y(x) = \alpha e^{5x}$, with $\alpha = \pm e^{C}$ an arbitrary non-zero constant --- and $\alpha = 0$ is the constant solution $y = 0$ dropped when dividing by $y$.
	This is the usual exponential growth (or decay, if the coefficient is negative) that appears in most of the models of Section 1.

	\paragraph{Second-Order Equations}
	\noindent\\
	The same idea sometimes works on a second-order equation in which $y$ itself does not appear, only its derivatives.
	Substituting
	$$
	z = \frac{dy}{dx}
	\qquad \text{so that} \qquad
	\frac{d^2y}{dx^2} = \frac{dz}{dx}
	$$
	turns it into a \textit{first-order} equation in $z$, which may be separable.
	Once $z$ is found, a last integration gives $y$ (and hence a second constant of integration).

	\section{(Homogeneous) Second-Order Linear Differential Equations}
	
	A second-order homogeneous linear differential equation is of the form:
	$$
	a\frac{d^2y}{dx^2} + b\frac{dy}{dx} + cy = 0
	$$
	The characteristic equation associated with this differential equation is:
	$$
	ar^2 + br + c = 0
	$$
	The nature of the solutions depends on the discriminant $\Delta = b^2 - 4ac$:
	\begin{itemize}
		\item If $\Delta > 0$, the roots $r_1$ and $r_2$ are real and distinct.
		The solution is: $y(x) = C_1e^{r_1x} + C_2e^{r_2x}$.
		\item If $\Delta = 0$, the roots are real and repeated $r_1 = r_2 = r$.
		The solution is: $y(x) = (C_1 + C_2x)e^{rx}$.
		\item If $\Delta < 0$, the roots are complex $r_{1,2} = \alpha \pm i\beta$.
		The solution is: $y(x) = e^{\alpha x}(C_1\cos(\beta x) + C_2\sin(\beta x))$.
	\end{itemize}

	\noindent \textbf{Example ($\Delta > 0$):}
	For $y'' - y' - 6y = 0$, the characteristic equation is $r^2 - r - 6 = 0$, with $\Delta = 1 + 24 = 25 > 0$, so
	$$
	r = \frac{1 \pm 5}{2} \quad \Longrightarrow \quad r_1 = 3 \ \text{ and } \ r_2 = -2
	$$
	and the general solution is
	$$
	y(x) = C_1 e^{3x} + C_2 e^{-2x}, \qquad C_1, C_2 \in \R
	$$

	\noindent \textbf{Example ($\Delta = 0$):}
	For $y'' - 6y' + 9y = 0$, the characteristic equation is $r^2 - 6r + 9 = 0$, i.e. $(r-3)^2 = 0$, with $\Delta = 36 - 36 = 0$, so $r = 3$ is a repeated root and
	$$
	y(x) = (C_1 + C_2 x) e^{3x}, \qquad C_1, C_2 \in \R
	$$
	Note the extra factor $x$: without it we would only get \textit{one} family of solutions, whereas a second-order equation always has \textit{two} independent ones.

	\noindent \textbf{Example ($\Delta < 0$):}
	For $y'' - 2y' + 10y = 0$, the characteristic equation is $r^2 - 2r + 10 = 0$, with $\Delta = 4 - 40 = -36 < 0$, so
	$$
	r = \frac{2 \pm \sqrt{-36}}{2} = \frac{2 \pm 6i}{2} = 1 \pm 3i
	\qquad \Longrightarrow \qquad \alpha = 1, \ \beta = 3
	$$
	and the general solution is
	$$
	y(x) = e^{x}\left(C_1 \cos(3x) + C_2 \sin(3x)\right), \qquad C_1, C_2 \in \R
	$$
	Complex roots always give oscillations: $\beta$ sets the frequency of the oscillation, and $\alpha$ says whether it is amplified ($\alpha > 0$) or damped ($\alpha < 0$).

	\section{Second-Order Linear Differential Equations}
	
	For the general second-order linear differential equation:
	$$
	a\frac{d^2y}{dx^2} + b\frac{dy}{dx} + cy = g(x)
	$$
	The solution is the sum of the complementary function $y_c(x)$
	(the solution to the associated homogeneous equation)
	and a particular solution $y_p(x)$:
	$$
	y(x) = y_c(x) + y_p(x)
	$$
	
	The method for finding $y_p(x)$ depends on the form of $g(x)$:
	\begin{itemize}
		\item If $g(x)$ is a polynomial, use the method of undetermined coefficients.
		\item If $g(x)$ is of an exponential or trigonometric form, use an answer of a similar form.
		\item For other forms of $g(x)$, the method of variation of parameters can be used.
	\end{itemize}

	\noindent \textbf{Example ($g$ polynomial):}
	Solve $y'' - y' - 6y = 12x + 2$.

	The associated homogeneous equation was solved above:
	$$
	y_c(x) = C_1 e^{3x} + C_2 e^{-2x}
	$$
	Since $g(x) = 12x + 2$ is a polynomial of degree 1, we look for a particular solution of the same form, $y_p(x) = ax + b$.
	Then $y_p' = a$, $y_p'' = 0$, and
	$$
	y_p'' - y_p' - 6y_p = -a - 6(ax + b) = -6a \, x + (-a - 6b)
	$$
	We want this to equal $12x + 2$ \textit{for every $x$}, so we identify the coefficients:
	$$
	\begin{cases}
		-6a = 12 \\
		-a - 6b = 2
	\end{cases}
	\qquad \Longrightarrow \qquad
	a = -2, \quad b = 0
	$$
	so $y_p(x) = -2x$, and the general solution is
	$$
	y(x) = y_c(x) + y_p(x) = C_1 e^{3x} + C_2 e^{-2x} - 2x
	$$

	\noindent \textbf{Example ($g$ exponential):}
	Solve $y'' - y' - 6y = 4e^{x}$.

	The complementary function $y_c$ is the same as above.
	Since $g(x) = 4e^x$ is an exponential, we look for $y_p(x) = A e^{x}$; then $y_p' = y_p'' = Ae^x$, and
	$$
	y_p'' - y_p' - 6y_p = (A - A - 6A) e^{x} = -6A \, e^{x} = 4 e^{x}
	\qquad \Longrightarrow \qquad
	A = -\frac{2}{3}
	$$
	so the general solution is
	$$
	y(x) = C_1 e^{3x} + C_2 e^{-2x} - \frac{2}{3} e^{x}
	$$
	\textit{Warning:} this works because $e^{x}$ is not itself a solution of the homogeneous equation.
	If $g(x)$ had been $4e^{3x}$, the guess $Ae^{3x}$ would give $0 = 4e^{3x}$, which is impossible; one then tries $y_p(x) = Ax e^{3x}$ instead (the same trick as for a repeated root).

	\paragraph{Initial Conditions}
	
	Given the general solution:
	$$
	y(x) = y_c(x) + y_p(x)
	$$
	And initial conditions $y(x_0) = y_0$ and $\frac{dy}{dx}\Big|_{x=x_0} = y'_0$, we determine the constants $C_1$ and $C_2$ in $y_c(x)$:
	$$
	y(x_0) = y_c(x_0) + y_p(x_0) = y_0
	$$
	$$
	\frac{dy}{dx}\Big|_{x=x_0} = y'_c(x_0) + y'_p(x_0) = y'_0
	$$
	Solving these two equations will yield the specific values of $C_1$ and $C_2$ that satisfy the initial conditions, leading to the particular solution for the problem.

	\noindent \textbf{Example:}
	Take the equation $y'' - y' - 6y = 12x + 2$ solved above, with the initial conditions $y(0) = 1$ and $y'(0) = 0$.

	The general solution and its derivative are
	$$
	y(x) = C_1 e^{3x} + C_2 e^{-2x} - 2x
	\qquad \text{and} \qquad
	y'(x) = 3C_1 e^{3x} - 2C_2 e^{-2x} - 2
	$$
	The initial conditions give
	$$
	\begin{cases}
		y(0) = C_1 + C_2 = 1 \\
		y'(0) = 3C_1 - 2C_2 - 2 = 0
	\end{cases}
	$$
	Substituting $C_2 = 1 - C_1$ in the second equation: $3C_1 - 2 + 2C_1 = 2$, i.e. $5C_1 = 4$, so
	$$
	C_1 = \frac{4}{5}, \qquad C_2 = \frac{1}{5}
	$$
	and the solution of the initial value problem is
	$$
	y(x) = \frac{4}{5} e^{3x} + \frac{1}{5} e^{-2x} - 2x
	$$
	\textit{Check:} $y(0) = \frac{4}{5} + \frac{1}{5} = 1$ \checkmark \ and $y'(0) = \frac{12}{5} - \frac{2}{5} - 2 = 2 - 2 = 0$. \checkmark

	\textit{Note:} the two constants are found \textit{at the very end}, once the particular solution $y_p$ has been added.
	Applying the initial conditions to $y_c$ alone is a common mistake, and gives a wrong answer.
	
	
\end{document}
